\begin{abstract}
Prediction Shift Backdoor Detection (PSBD) flags a poisoned input by how little its
prediction moves when the network is perturbed at inference. It was built for ResNets,
where the residual stream is the one obvious place to put that perturbation. A
transformer block offers many places and several kinds of perturbation, and we show that
this choice decides whether the detector works. Masking whole tokens at the input of
every attention block reaches a mean AUROC of \HeadlineAurocAdaptive{} over
\PanelCellsCached{} backdoored ViT-B/16 models on \PanelDatasetsCached{} datasets and
poison rates of \PanelRates{}, against \PublishedAurocAdaptive{} for the placement the
original work used. The paired gain of \HeadlineGainAdaptiveAuroc{}
[\HeadlineGainAdaptiveAurocLow{}, \HeadlineGainAdaptiveAurocHigh{}] grows to
\GainsLowestRate{} at \PanelRateLow{} poisoning, the regime a defender most needs, and
the placement transfers to Swin-S, where it reaches \SwinRecommendedAurocAdaptive{} over
\SwinCells{} models. Against \DetectorsComparedCount{} input-level detectors ported to
the same models and splits it ranks \DetectorsAurocRankOurs{} of \DetectorsAurocDefensesRanked{} defenses
by AUROC and by the true-positive rate at both false-positive budgets we report. We
explain why by following the backdoor through the network. A patch trigger's evidence
sits in its own tokens of the residual stream until the last blocks. Token masking at the
attention input leaves that stream untouched and changes only
\SurvivalTmBadnetATwooTriggered{} of triggered BadNets predictions while it changes
\SurvivalTmClean{} of clean ones, whereas dropout on the stream changes
\SurvivalRdBadnetATwooTriggered{} of them. A
union of probes keeps detection at \AdaptiveVitUnion{} AUROC against an attacker trained
to evade a single probe.
\end{abstract}
